\subsubsection{毛管波制約の導出：分散関係から CFL 条件へ}
\label{app:capillary_cfl}
% ═══════════════════════════════════════════════════════════════════════════════

\textbf{Step 1：毛管波の分散関係}

気液界面上の微小振幅波（毛管波）の線形安定解析から，
角周波数 $\omega$ と波数 $k$ の関係（分散関係）は：
\begin{equation}
  \omega^2 = \frac{\sigma\,k^3}{\rho_l + \rho_g}
  \label{eq:capillary_dispersion_app}
\end{equation}
位相速度 $c_\sigma = \omega/k = \sqrt{\sigma k/(\rho_l+\rho_g)}$，
群速度 $c_g = d\omega/dk = (3/2)c_\sigma$．

\textbf{Step 2：格子上の最大分解可能波数}

格子間隔 $\Delta x$ の離散格子で分解できる最高波数はナイキスト波数
$k_{\max} = \pi/\Delta x$ である．

\textbf{Step 3：CFL 条件の適用}

$k_{\max}$ に対応する群速度で CFL 条件 $c_g(k_{\max})\,\Delta t \le \Delta x$ を適用する：
\begin{align*}
  c_g(k_{\max}) &= \frac{3}{2}\sqrt{\frac{\sigma\pi/\Delta x}{\rho_l+\rho_g}} \\
  \Delta t &\le \frac{\Delta x}{c_g(k_{\max})}
  = \frac{2}{3}\sqrt{\frac{(\rho_l+\rho_g)\,\Delta x^3}{\sigma\pi}}
\end{align*}
厳密な係数は $2/(3\sqrt{\pi}) \approx 0.376$ であるが，
文献で広く用いられる簡略形 $1/\sqrt{2\pi} \approx 0.399$ を採用すると
式 \eqref{eq:dt_sigma} を得る（$6\%$ 大きい $\Delta t$ を許すが，
毛管波 CFL の安全係数 $C_\sigma \in (0,1]$ を
$\Delta t \le C_\sigma\Delta t_\sigma$ として掛けるため実用上問題ない）：
\[
  \Delta t_\sigma = \sqrt{\frac{(\rho_l+\rho_g)\,\Delta x^3}{2\pi\sigma}}
\]

\textbf{物理的解釈：} $\Delta t_\sigma \propto \Delta x^{3/2}$ は分散関係 $\omega \propto k^{3/2}$ に由来する．
格子を半分にすると $\Delta t_\sigma$ は $1/2^{3/2} \approx 0.354$ 倍（$\approx 2.8$ 倍制約強化）．

\medskip
\textbf{大密度比（$\rho_l \gg \rho_g$）での解釈：}
$\rho_l/\rho_g = 1000$（水/空気）では $\rho_l + \rho_g \approx \rho_l$（気相寄与は $0.1\%$ 以下）．
表面張力が\textbf{復元力}（ばね），液相 $\rho_l$ が\textbf{慣性}（質量）を担い，
気相の慣性は無視できる．密度比 1000 の系では密度比 1 の系より
$\sqrt{1000} \approx 32$ 倍大きな $\Delta t_\sigma$ が許される（制約緩和）．\qed

% ═══════════════════════════════════════════════════════════════════════════════
%% 旧 appendix_numerics_schemes_s4.tex（ALE 効果）を統合
% ═══════════════════════════════════════════════════════════════════════════════

\subsubsection{時間変化格子と ALE 効果}
\label{app:grid_ale}

本付録は §\ref{sec:grid_gen} の格子が時間発展する場合の ALE（Arbitrary Lagrangian-Eulerian）補正を説明する．
\textbf{本稿の標準経路では，界面適合格子を毎ステップ再生成し，保存的再写像と速度再投影で閉じる}．
時間変化格子では，§\ref{sec:grid_ale} の
離散幾何保存則，ALE 保存更新または同等の保存的再写像，全変数再写像，
および界面面データ再構築ゲートを全て満たす必要がある．

ALE 定式化では格子速度 $\bm{v}_\text{mesh}$ を輸送方程式に追加する必要がある：
$\partial\psi/\partial t + \nabla\cdot((\bm{u} - \bm{v}_\text{mesh})\psi) = 0$．
同じ相対速度 $\bm{u}-\bm{v}_\text{mesh}$ は運動量・密度輸送にも現れるため，
$\psi$ だけを補間する格子追従は閉じた保存系ではない．

\textbf{誤差の発生機構：}
1ステップあたりの界面移動量を $\delta = |\bm{u}|\Delta t$ とすると，
ALE 補正を省略することで $\bm{v}_\text{mesh}$ による追加対流量が無視される：
\[
  \varepsilon_{\text{ALE}} \sim |\bm{v}_\text{mesh}| \cdot |\bnabla\psi| \cdot \Delta t \sim U_\Gamma \cdot \frac{1}{\varepsilon} \cdot \Delta t
\]
ここで $|\bnabla\psi| \sim 1/(4\varepsilon)$（界面近傍プロファイル），$U_\Gamma$ は界面速度である．

\textbf{近似が有効な条件：}
\begin{center}
\PaperTableSetup
\begin{tabular}{>{\raggedright\arraybackslash}p{48mm}c>{\raggedright\arraybackslash}p{42mm}}
\hline
条件 & 移動量 $\delta/h$ & ALE 省略の妥当性 \\
\hline
緩やかな界面変形 & $< 0.1$ & 良好（誤差 $< 10\%$ of $\Delta t^2$ 項） \\
標準的な問題 & $\approx 0.5$ & 条件付きで可（格子収束確認要） \\
大密度比 ($\rho_l/\rho_g{=}1000$) の高速界面 & $> 1$ & 精度劣化の可能性大（ALE 実装推奨） \\
\hline
\end{tabular}
\end{center}

$\rho_l/\rho_g = 1000$ の場合，液相変形が速く $\delta/h > 1$ となることがある．
この場合，単純な $\psi$ 補間だけの格子更新は実質的に $O(\Delta t)$ 精度（1次）となりうるため，
標準経路では体積重み付き質量補正，速度再投影，履歴リセット，面幾何再構築を同時に行う．

\textbf{Mode 1（時間固定格子，比較経路）：}
格子速度 $\bm{v}_\text{mesh} = 0$ として格子を更新しない．ALE 誤差はゼロであり，全体時間精度 $O(\Delta t^2)$ が保持される（§\ref{sec:time_accuracy_table}）．
ただし，格子は初期水準集合値 $\phi^{(0)}$ から生成されるため，界面が初期精細化領域から大きく離れる問題（例：上昇気泡ベンチマーク）では経時的に空間解像度が失われる．
長時間・長距離移動を伴う計算では，Mode 2 の界面追随格子を標準経路とする．

\textbf{Mode 2（タイムステップごとに格子を再生成，標準経路）：}
上表の $\delta/h > 1$ 条件に該当する場合，$O(\Delta t)$ に精度が劣化し，本稿の全体時間精度 $O(\Delta t^2)$ の主張（§\ref{sec:time_accuracy_table}）と矛盾する．
$\rho_l/\rho_g = 1000$（上昇気泡系）では特に注意が必要であり，Mode 2 では ALE 補正または保存的再写像を必須とする．
Mode 2 は $\psi$ 転写だけでは閉じない．
$u,v,p,\delta p,\rho,\mu,\kappa_{lg},\hat{\bm n}_{lg},\phi,\psi$ の保存的リマップと，
HFE 延長文脈・jump face cache の現在格子上での再構築を同時に満たす必要がある．
これらが欠ける場合，格子追従の利得より operator mismatch が支配的になりうるため，
式~\eqref{eq:ale_grid_gcl}--\eqref{eq:regrid_face_geometry} を標準経路の成立条件とする．

% ═══════════════════════════════════════════════════════════════════════════════
